\documentclass[12pt,a4paper]{report}

% ============================== PACKAGES ==============================
\usepackage[utf8]{inputenc}
\usepackage[T1]{fontenc}
\usepackage[french]{babel}
\usepackage{geometry}
\geometry{margin=2.5cm}
\usepackage{graphicx}
\usepackage{float}
\usepackage{amsmath,amssymb}
\usepackage{booktabs}
\usepackage{array}
\usepackage{listings}
\usepackage{xcolor}
\usepackage{hyperref}
\usepackage{caption}
\usepackage{subcaption}
\usepackage{tikz}
\usetikzlibrary{shapes,arrows,positioning,fit,calc}
\usepackage{fancyhdr}
\usepackage{titlesec}


% ============================== LISTINGS STYLE ==============================
\lstdefinestyle{vhdl}{
  language=VHDL,
  basicstyle=\ttfamily\footnotesize,
  keywordstyle=\color{blue}\bfseries,
  commentstyle=\color{green!50!black}\itshape,
  stringstyle=\color{red},
  numbers=left,
  numberstyle=\tiny\color{gray},
  numbersep=5pt,
  frame=single,
  breaklines=true,
  tabsize=2,
  showstringspaces=false,
  captionpos=b
}

% ============================== HEADER/FOOTER ==============================
\pagestyle{fancy}
\fancyhf{}
\fancyhead[L]{Co-processeur Cryptographique RSA}
\fancyhead[R]{RISC-V + Crypto}
\fancyfoot[C]{\thepage}

% ============================== TITLE PAGE ==============================
\begin{document}


\begin{titlepage}
\centering
\vspace*{2cm}

{\Huge\bfseries Rapport de Projet\\[0.5cm]}
{\LARGE Co-processeur Cryptographique RSA\\pour Architecture RISC-V\\[2cm]}

{\Large
\begin{tabular}{rl}
\textbf{Module :}    & Projet de Fin d'Études \\
\textbf{Filière :}   & Électronique / Systèmes Embarqués \\
\textbf{Encadrant :} & \textit{[Nom de l'encadrant]} \\
\textbf{Réalisé par :} & \textit{[Votre nom]} \\
\textbf{Année :}     & 2024--2025 \\
\end{tabular}
}

\vfill
{\large\textit{Université / École}}
\end{titlepage}

% ============================== TOC ==============================
\tableofcontents
\newpage


% ======================================================================
%                        CHAPTER 1 : INTRODUCTION
% ======================================================================
\chapter{Introduction}

\section{Contexte du projet}
Ce projet s'inscrit dans le cadre de la conception d'un co-processeur
cryptographique dédié à l'algorithme RSA, destiné à être intégré à un
processeur RISC-V. L'objectif est de décharger le processeur principal des
opérations d'exponentiation modulaire coûteuses en cycles, permettant ainsi
un chiffrement/déchiffrement RSA accéléré par matériel.

\section{Objectifs}
\begin{itemize}
  \item Implémenter l'algorithme RSA complet en VHDL (génération de clés,
        chiffrement, déchiffrement).
  \item Utiliser la multiplication de Montgomery pour l'exponentiation
        modulaire efficace.
  \item Concevoir une architecture modulaire et paramétrable (largeur de
        clé configurable).
  \item Valider le fonctionnement par simulation (testbench VHDL).
  \item Préparer l'intégration avec un bus AXI/Wishbone vers le RISC-V.
\end{itemize}

\section{Organisation du rapport}
Le chapitre~\ref{ch:theory} présente les fondements mathématiques du RSA
et de Montgomery. Le chapitre~\ref{ch:archi} détaille l'architecture
globale. Les chapitres~\ref{ch:entities} à~\ref{ch:top} décrivent chaque
entité VHDL. Le chapitre~\ref{ch:sim} présente les résultats de simulation.
Enfin, le chapitre~\ref{ch:conclusion} conclut et propose des perspectives.


% ======================================================================
%                   CHAPTER 2 : FONDEMENTS THÉORIQUES
% ======================================================================
\chapter{Fondements Théoriques}
\label{ch:theory}

\section{Algorithme RSA}
L'algorithme RSA repose sur la difficulté de factoriser le produit de
deux grands nombres premiers.

\subsection{Génération de clés}
\begin{enumerate}
  \item Choisir deux nombres premiers $p$ et $q$ de $k/2$ bits chacun.
  \item Calculer $N = p \times q$ \quad (module RSA, $k$ bits).
  \item Calculer $\varphi(N) = (p-1)(q-1)$.
  \item Choisir $e$ tel que $\gcd(e, \varphi(N)) = 1$ (exposant public,
        typiquement $e = 65537$).
  \item Calculer $d = e^{-1} \bmod \varphi(N)$ (exposant privé, via
        l'algorithme d'Euclide étendu).
\end{enumerate}
La clé publique est $(N, e)$, la clé privée est $(N, d)$.

\subsection{Chiffrement et déchiffrement}
\begin{align}
  \text{Chiffrement :} \quad c &= m^e \bmod N \\
  \text{Déchiffrement :} \quad m &= c^d \bmod N
\end{align}
où $m < N$ est le message en clair et $c$ est le chiffré.


\section{Multiplication de Montgomery}

La multiplication de Montgomery permet d'effectuer des réductions modulaires
sans division explicite, en travaillant dans un domaine transformé.

\subsection{Principe}
Soit $R = 2^K$ avec $\gcd(R, N) = 1$ (N impair). La forme de Montgomery
d'un entier $a$ est $\bar{a} = a \cdot R \bmod N$.

La \textbf{réduction de Montgomery} calcule :
\[
  \text{MonPro}(a, b) = a \cdot b \cdot R^{-1} \bmod N
\]

\subsection{Algorithme bit-sériel}
L'implémentation utilise l'algorithme bit-sériel de Montgomery :
\begin{enumerate}
  \item $S \leftarrow 0$
  \item Pour $i = 0$ à $K-1$ :
  \begin{enumerate}
    \item Si $X_i = 1$ : $S \leftarrow S + Y$
    \item Si $S_0 = 1$ : $S \leftarrow S + M$
    \item $S \leftarrow S \gg 1$
  \end{enumerate}
  \item Si $S \geq M$ : $S \leftarrow S - M$
  \item Retourner $S$
\end{enumerate}

Cet algorithme traite un bit par cycle d'horloge, nécessitant $K+2$
cycles par multiplication (K cycles de calcul + 1 cycle de correction
+ 1 cycle de sortie).


\section{Exponentiation modulaire par Montgomery}

L'exponentiation $X^E \bmod N$ utilise la méthode \textit{square-and-multiply}
(LSB-first dans notre implémentation) avec deux instances de Montgomery :

\begin{enumerate}
  \item Précalculer $R^2 \bmod N$ par doublement itéré ($2K$ cycles).
  \item $\bar{Z} = \text{MonPro}(1, R^2) = R \bmod N$ (forme de Montgomery de 1).
  \item $\bar{S} = \text{MonPro}(X, R^2) = X \cdot R \bmod N$ (forme de
        Montgomery de X).
  \item Pour chaque bit $i$ de $E$ (de 0 à $K_{exp}-1$) :
  \begin{itemize}
    \item $\bar{S} \leftarrow \text{MonPro}(\bar{S}, \bar{S})$ \quad (carré)
    \item Si $E_i = 1$ : $\bar{Z} \leftarrow \text{MonPro}(\bar{Z}, \bar{S})$
          \quad (multiplication)
  \end{itemize}
  \item $Z = \text{MonPro}(\bar{Z}, 1)$ \quad (retour au domaine normal).
\end{enumerate}

\section{Test de primalité de Miller-Rabin}

Pour la génération de clés, les candidats premiers sont testés par
l'algorithme de Miller-Rabin avec $t$ témoins ($a = 2, 3, 5, 7$).

Soit $n - 1 = 2^s \cdot d$ avec $d$ impair. Pour chaque témoin $a$ :
\begin{enumerate}
  \item $x = a^d \bmod n$
  \item Si $x = 1$ ou $x = n-1$ : passer au témoin suivant.
  \item Pour $r = 1$ à $s-1$ : $x = x^2 \bmod n$. Si $x = n-1$ : passer.
  \item Si aucun passage : $n$ est composé.
\end{enumerate}

\section{Algorithme d'Euclide étendu (EXT\_GCD)}
Calcule $d = e^{-1} \bmod \varphi(N)$ par l'algorithme itératif classique
avec division et multiplication séquentielles (un bit par cycle).


% ======================================================================
%                   CHAPTER 3 : ARCHITECTURE GLOBALE
% ======================================================================
\chapter{Architecture Globale}
\label{ch:archi}

\section{Vue d'ensemble}

L'architecture est hiérarchique et modulaire. Le schéma ci-dessous montre
les blocs principaux et leurs interconnexions :

\begin{figure}[H]
\centering
\begin{tikzpicture}[
  block/.style={draw, thick, minimum width=3.5cm, minimum height=1.2cm,
                align=center, rounded corners=3pt},
  arr/.style={->, thick, >=stealth}
]

% RSA_TOP
\node[block, fill=blue!10, minimum width=12cm, minimum height=8cm]
  (top) at (0,0) {};
\node[above] at (top.north) {\textbf{\large RSA\_TOP}};

% RSA_KEYGEN block
\node[block, fill=orange!20] (keygen) at (-3, 2) {RSA\_KEYGEN};

% Sub-blocks of KEYGEN
\node[block, fill=yellow!20, font=\small] (rng) at (-5, 0) {RANDOM\_GEN\\(LFSR 128b)};
\node[block, fill=yellow!20, font=\small] (mr) at (-3, -1.5) {MILLER\_RABIN};
\node[block, fill=yellow!20, font=\small] (gcd) at (-1, 0) {EXT\_GCD};

% RSA core
\node[block, fill=green!15] (rsa) at (3, 2) {RSA};

% Montgomery EXP
\node[block, fill=green!30, font=\small] (exp) at (3, 0) {MOD\_MONTGOMERY\\\_EXP};

% Montgomery MULT
\node[block, fill=red!15, font=\small] (mult1) at (2, -2.2) {MONTGOMERY\\\_MULT (×)};
\node[block, fill=red!15, font=\small] (mult2) at (4.5, -2.2) {MONTGOMERY\\\_MULT (²)};

% Arrows
\draw[arr] (keygen) -- node[above]{N, e, d} (rsa);
\draw[arr] (rng) -- (keygen);
\draw[arr] (mr) -- (keygen);
\draw[arr] (gcd) -- (keygen);
\draw[arr] (rsa) -- (exp);
\draw[arr] (exp) -- (mult1);
\draw[arr] (exp) -- (mult2);

% External I/O
\draw[arr] (-6.5, 2) -- node[above]{\small start, seed} (keygen.west);
\draw[arr] (-6.5, 3) -- node[above]{\small i\_message, mode, start} (top.west |- rsa);
\draw[arr] (rsa.east) -- ++(1,0) node[right]{\small o\_result, o\_done};

\end{tikzpicture}
\caption{Architecture globale du co-processeur RSA}
\label{fig:archi_globale}
\end{figure}


\section{Hiérarchie des entités}

\begin{table}[H]
\centering
\caption{Hiérarchie des fichiers VHDL}
\label{tab:hierarchy}
\begin{tabular}{lll}
\toprule
\textbf{Fichier} & \textbf{Entité} & \textbf{Rôle} \\
\midrule
\texttt{RSA\_PKG.vhd}       & (package)           & Constantes, types, fonctions \\
\texttt{RSA\_TOP.vhd}       & \texttt{RSA\_TOP}   & Top-level (keygen + RSA) \\
\texttt{RSA\_KEYGEN.vhd}    & \texttt{RSA\_KEYGEN} & Génération de clés complète \\
\texttt{RANDOM\_GEN.vhd}    & \texttt{RANDOM\_GEN} & PRNG (LFSR 128 bits) \\
\texttt{MILLER\_RABIN.vhd}  & \texttt{MILLER\_RABIN} & Test de primalité \\
\texttt{EXT\_GCD.vhd}       & \texttt{EXT\_GCD}   & Euclide étendu ($d = e^{-1}$) \\
\texttt{RSA.vhd}            & \texttt{RSA}        & Wrapper exponentiation mod \\
\texttt{MOD\_MONTGOMERY\_EXP.vhd} & \texttt{MOD\_MONTGOMERY\_EXP} & Exponentiation Montgomery \\
\texttt{MONTGOMERY\_MULT.vhd} & \texttt{MONTGOMERY\_MULT} & Multiplieur bit-sériel \\
\texttt{TB\_RSA\_TOP.vhd}   & \texttt{TB\_RSA\_TOP} & Testbench principal \\
\bottomrule
\end{tabular}
\end{table}

\section{Paramètres de conception}

\begin{table}[H]
\centering
\caption{Paramètres configurables}
\label{tab:params}
\begin{tabular}{lcp{7cm}}
\toprule
\textbf{Paramètre} & \textbf{Valeur} & \textbf{Description} \\
\midrule
\texttt{PRIME\_WIDTH}   & 512   & Largeur des premiers $p$, $q$ (bits) \\
\texttt{KEY\_WIDTH}     & 1024  & Largeur du module $N = p \times q$ \\
\texttt{NUM\_WITNESSES} & 4     & Nombre de témoins Miller-Rabin \\
\texttt{e}              & 65537 & Exposant public fixe \\
LFSR                    & 128 bits & Graine PRNG \\
\bottomrule
\end{tabular}
\end{table}


% ======================================================================
%              CHAPTER 4 : DESCRIPTION DES ENTITÉS
% ======================================================================
\chapter{Description Détaillée des Entités}
\label{ch:entities}

% ----------------------------------------------------------------------
\section{RSA\_PKG --- Package de constantes}
\label{sec:pkg}

Ce package centralise les constantes globales et la fonction
\texttt{mod\_inverse\_fn} (Euclide étendu) utilisée à l'élaboration.

\begin{table}[H]
\centering
\caption{Constantes définies dans \texttt{RSA\_PKG}}
\begin{tabular}{lcl}
\toprule
\textbf{Nom} & \textbf{Valeur} & \textbf{Usage} \\
\midrule
\texttt{PRIME\_WIDTH} & 512  & Largeur des premiers \\
\texttt{KEY\_WIDTH}   & 1024 & $= 2 \times$ \texttt{PRIME\_WIDTH} \\
\bottomrule
\end{tabular}
\end{table}

% ----------------------------------------------------------------------
\section{RANDOM\_GEN --- Générateur pseudo-aléatoire}
\label{sec:rng}

\subsection{Fonction}
Génère des nombres pseudo-aléatoires de \texttt{WIDTH} bits à partir
d'un LFSR (Linear Feedback Shift Register) de 128 bits. Le polynôme
de rétroaction est :
\[
  x^{128} + x^{126} + x^{101} + x^{99} + 1
\]
La sortie est forcée impaire (LSB $= 1$) et de pleine longueur (MSB $= 1$).

\subsection{Interface}
\begin{table}[H]
\centering
\caption{Ports de \texttt{RANDOM\_GEN}}
\begin{tabular}{llcp{5cm}}
\toprule
\textbf{Port} & \textbf{Direction} & \textbf{Largeur} & \textbf{Description} \\
\midrule
\texttt{clk}    & in  & 1   & Horloge \\
\texttt{rst}    & in  & 1   & Reset asynchrone \\
\texttt{seed}   & in  & 128 & Graine initiale (non nulle) \\
\texttt{load}   & in  & 1   & Chargement de la graine \\
\texttt{start}  & in  & 1   & Demande d'un nouveau nombre \\
\texttt{o\_done} & out & 1   & Sortie valide \\
\texttt{o\_rng}  & out & WIDTH & Nombre aléatoire généré \\
\bottomrule
\end{tabular}
\end{table}


\subsection{Machine d'états}
\begin{enumerate}
  \item \textbf{IDLE} : attente d'une requête \texttt{start}.
  \item \textbf{STEP} : avance le LFSR et accumule 128 bits par cycle
        ($\lceil\text{WIDTH}/128\rceil$ cycles au total).
  \item \textbf{FINALIZE} : force MSB $= 1$ et LSB $= 1$.
  \item \textbf{DONE\_ST} : présente la sortie pendant un cycle.
\end{enumerate}

\subsection{Simulation}
\begin{figure}[H]
\centering
% === INSÉRER ICI LA CAPTURE DE SIMULATION ===
\fbox{\parbox{0.9\textwidth}{\centering\vspace{3cm}
  \textit{[Capture chronogramme RANDOM\_GEN : signaux clk, start, o\_done, o\_rng]}
\vspace{3cm}}}
\caption{Chronogramme de simulation de \texttt{RANDOM\_GEN}}
\label{fig:sim_rng}
\end{figure}

% ----------------------------------------------------------------------
\section{MILLER\_RABIN --- Test de primalité}
\label{sec:mr}

\subsection{Fonction}
Teste si un nombre $n$ de $K$ bits est probablement premier en utilisant
l'algorithme de Miller-Rabin avec \texttt{NUM\_WITNESSES} témoins
fixes ($a \in \{2, 3, 5, 7\}$).

\subsection{Interface}
\begin{table}[H]
\centering
\caption{Ports de \texttt{MILLER\_RABIN}}
\begin{tabular}{llcp{5cm}}
\toprule
\textbf{Port} & \textbf{Dir.} & \textbf{Largeur} & \textbf{Description} \\
\midrule
\texttt{clk}     & in  & 1 & Horloge \\
\texttt{rst}     & in  & 1 & Reset asynchrone \\
\texttt{start}   & in  & 1 & Lancement du test \\
\texttt{i\_n}    & in  & K & Nombre à tester \\
\texttt{o\_done}  & out & 1 & Résultat disponible \\
\texttt{o\_prime} & out & 1 & \texttt{'1'} = prob. premier \\
\bottomrule
\end{tabular}
\end{table}


\subsection{Machine d'états}
\begin{enumerate}
  \item \textbf{IDLE} : attente.
  \item \textbf{CHECK\_TRIVIAL} : cas triviaux ($n=2,3$, $n$ pair).
  \item \textbf{DECOMPOSE} : décomposition $n-1 = 2^s \cdot d$ (un shift/cycle).
  \item \textbf{LOAD\_WITNESS} : charge le témoin $a_i$.
  \item \textbf{ST\_EXP\_LAUNCH/WAIT} : calcule $x = a^d \bmod n$ via
        \texttt{MOD\_MONTGOMERY\_EXP}.
  \item \textbf{CHECK\_INITIAL} : vérifie $x = 1$ ou $x = n-1$.
  \item \textbf{ST\_SQ\_LAUNCH/WAIT} : élévations au carré successives
        $x = x^2 \bmod n$.
  \item \textbf{CHECK\_SQUARE} : vérifie $x = n-1$ après chaque carré.
  \item \textbf{COMPOSITE / PROBABLY\_PRIME} : verdict final.
  \item \textbf{DONE\_ST} : sortie du résultat.
\end{enumerate}

\subsection{Simulation}
\begin{figure}[H]
\centering
\fbox{\parbox{0.9\textwidth}{\centering\vspace{3cm}
  \textit{[Capture chronogramme MILLER\_RABIN : signaux clk, start, i\_n,
  o\_done, o\_prime, state]}
\vspace{3cm}}}
\caption{Chronogramme de simulation de \texttt{MILLER\_RABIN}}
\label{fig:sim_mr}
\end{figure}

% ----------------------------------------------------------------------
\section{EXT\_GCD --- Euclide étendu}
\label{sec:gcd}

\subsection{Fonction}
Calcule l'inverse modulaire $d = e^{-1} \bmod \varphi(N)$ en utilisant
l'algorithme d'Euclide étendu itératif. La division et la multiplication
internes sont réalisées en séquentiel (un bit par cycle).

\subsection{Interface}
\begin{table}[H]
\centering
\caption{Ports de \texttt{EXT\_GCD}}
\begin{tabular}{llcp{5cm}}
\toprule
\textbf{Port} & \textbf{Dir.} & \textbf{Largeur} & \textbf{Description} \\
\midrule
\texttt{clk}     & in  & 1 & Horloge \\
\texttt{rst}     & in  & 1 & Reset asynchrone \\
\texttt{start}   & in  & 1 & Lancement \\
\texttt{i\_e}    & in  & K & Valeur à inverser ($e$) \\
\texttt{i\_phi}  & in  & K & Module ($\varphi(N)$) \\
\texttt{o\_done}  & out & 1 & Résultat prêt \\
\texttt{o\_valid} & out & 1 & \texttt{'1'} si inverse existe \\
\texttt{o\_d}     & out & K & Résultat $e^{-1} \bmod \varphi$ \\
\bottomrule
\end{tabular}
\end{table}


\subsection{Machine d'états}
\begin{enumerate}
  \item \textbf{IDLE} : attente de \texttt{start}.
  \item \textbf{CHECK\_ZERO} : si $r = 0$, l'algorithme est terminé.
  \item \textbf{DIV\_STEP} : division séquentielle (shift-and-subtract),
        un bit par cycle ($K$ cycles).
  \item \textbf{MUL\_INIT / MUL\_STEP} : multiplication séquentielle
        (shift-and-add) du quotient par $s$, un bit par cycle ($K$ cycles).
  \item \textbf{UPDATE} : mise à jour des registres EEA.
  \item \textbf{FINAL\_ADJUST} : rend le résultat positif ($d = d + \varphi$
        si $d < 0$).
  \item \textbf{DONE\_ST} : sortie.
\end{enumerate}

\subsection{Simulation}
\begin{figure}[H]
\centering
\fbox{\parbox{0.9\textwidth}{\centering\vspace{3cm}
  \textit{[Capture chronogramme EXT\_GCD : signaux clk, start, i\_e,
  i\_phi, o\_done, o\_valid, o\_d]}
\vspace{3cm}}}
\caption{Chronogramme de simulation de \texttt{EXT\_GCD}}
\label{fig:sim_gcd}
\end{figure}

% ----------------------------------------------------------------------
\section{MONTGOMERY\_MULT --- Multiplieur de Montgomery}
\label{sec:mont_mult}

\subsection{Fonction}
Calcule le produit de Montgomery :
\[
  Z = X \cdot Y \cdot R^{-1} \bmod M
\]
où $R = 2^K$. L'algorithme est bit-sériel : un bit de $X$ est traité par
cycle d'horloge.

\subsection{Interface}
\begin{table}[H]
\centering
\caption{Ports de \texttt{MONTGOMERY\_MULT}}
\begin{tabular}{llcp{5cm}}
\toprule
\textbf{Port} & \textbf{Dir.} & \textbf{Largeur} & \textbf{Description} \\
\midrule
\texttt{clk}    & in  & 1 & Horloge \\
\texttt{rst}    & in  & 1 & Reset asynchrone \\
\texttt{start}  & in  & 1 & Lancement \\
\texttt{i\_X}   & in  & K & Premier opérande \\
\texttt{i\_Y}   & in  & K & Second opérande \\
\texttt{i\_M}   & in  & K & Module (impair) \\
\texttt{o\_done} & out & 1 & Résultat prêt \\
\texttt{o\_Z}    & out & K & $X \cdot Y \cdot R^{-1} \bmod M$ \\
\bottomrule
\end{tabular}
\end{table}


\subsection{Machine d'états}
\begin{enumerate}
  \item \textbf{IDLE} : capture des entrées sur \texttt{start}.
  \item \textbf{RUN} : boucle de $K$ cycles (un bit de $X$ par cycle).
        À chaque itération :
        $S \leftarrow S + X_i \cdot Y$, puis
        si $S_0 = 1$ : $S \leftarrow S + M$, puis
        $S \leftarrow S \gg 1$.
  \item \textbf{DONE} : correction finale ($S \geq M \Rightarrow S - M$),
        puis \texttt{o\_done} = 1.
\end{enumerate}

\subsection{Latence}
$K + 2$ cycles par multiplication (K itérations + DONE + retour IDLE).

\subsection{Simulation}
\begin{figure}[H]
\centering
\fbox{\parbox{0.9\textwidth}{\centering\vspace{3cm}
  \textit{[Capture chronogramme MONTGOMERY\_MULT : signaux clk, start,
  i\_X, i\_Y, i\_M, o\_done, o\_Z, state]}
\vspace{3cm}}}
\caption{Chronogramme de simulation de \texttt{MONTGOMERY\_MULT}}
\label{fig:sim_mont}
\end{figure}

% ----------------------------------------------------------------------
\section{MOD\_MONTGOMERY\_EXP --- Exponentiation modulaire}
\label{sec:mod_exp}

\subsection{Fonction}
Calcule $Z = X^E \bmod M$ en utilisant la méthode \textit{square-and-multiply}
(LSB-first) avec deux instances de \texttt{MONTGOMERY\_MULT} (une pour le
carré, une pour la multiplication conditionnelle).

\subsection{Interface}
\begin{table}[H]
\centering
\caption{Ports de \texttt{MOD\_MONTGOMERY\_EXP}}
\begin{tabular}{llcp{5cm}}
\toprule
\textbf{Port} & \textbf{Dir.} & \textbf{Largeur} & \textbf{Description} \\
\midrule
\texttt{clk}    & in  & 1 & Horloge \\
\texttt{rst}    & in  & 1 & Reset asynchrone \\
\texttt{start}  & in  & 1 & Lancement \\
\texttt{i\_X}   & in  & K & Base \\
\texttt{i\_e}   & in  & K & Exposant \\
\texttt{i\_Mod} & in  & K & Module (impair) \\
\texttt{o\_done} & out & 1 & Résultat prêt \\
\texttt{o\_Z}    & out & K & $X^E \bmod M$ \\
\bottomrule
\end{tabular}
\end{table}


\subsection{Machine d'états}
\begin{enumerate}
  \item \textbf{IDLE} : capture de $X$, $E$, $M$.
  \item \textbf{PRECOMPUTE\_R2} : calcul de $R^2 \bmod M$ par $2K$
        doublements ($2K$ cycles).
  \item \textbf{INIT\_MONTG} : lance en parallèle
        $\bar{Z} = \text{MonPro}(1, R^2)$ et
        $\bar{S} = \text{MonPro}(X, R^2)$.
  \item \textbf{WAIT\_INIT} : attente des deux résultats ($K+2$ cycles).
  \item \textbf{STEP} : pour chaque bit de $E$, lance le carré
        $\bar{S} = \text{MonPro}(\bar{S}, \bar{S})$ et conditionnellement
        $\bar{Z} = \text{MonPro}(\bar{Z}, \bar{S})$.
  \item \textbf{WAIT\_STEP} : attente ($K+2$ cycles par itération).
  \item \textbf{FINAL} : $Z = \text{MonPro}(\bar{Z}, 1)$.
  \item \textbf{WAIT\_FINAL} : attente du résultat final.
  \item \textbf{DONE} : sortie de $o\_Z$.
\end{enumerate}

\subsection{Nombre de cycles}
\begin{table}[H]
\centering
\caption{Décomposition du nombre de cycles (\texttt{MOD\_MONTGOMERY\_EXP})}
\begin{tabular}{lr}
\toprule
\textbf{Phase} & \textbf{Cycles (K=1024)} \\
\midrule
PRECOMPUTE\_R2          & 2048 \\
INIT (2 × MonMul)      & 1026 \\
Boucle ($K$ itérations) & $K \times (K+2) + \text{popcount}(E) \times (K+2)$ \\
FINAL (1 × MonMul)     & 1026 \\
\midrule
\textbf{Total (pire cas)} & $\approx 2K^2 + 5K$ \\
\bottomrule
\end{tabular}
\end{table}

\subsection{Simulation}
\begin{figure}[H]
\centering
\fbox{\parbox{0.9\textwidth}{\centering\vspace{3cm}
  \textit{[Capture chronogramme MOD\_MONTGOMERY\_EXP : signaux clk, start,
  i\_X, i\_e, i\_Mod, o\_done, o\_Z, STATE]}
\vspace{3cm}}}
\caption{Chronogramme de simulation de \texttt{MOD\_MONTGOMERY\_EXP}}
\label{fig:sim_exp}
\end{figure}


% ----------------------------------------------------------------------
\section{RSA --- Wrapper d'exponentiation}
\label{sec:rsa}

\subsection{Fonction}
Entité enveloppe qui effectue $o\_result = i\_message^{i\_exp} \bmod i\_N$.
Elle réduit d'abord le message ($m \bmod N$) puis lance
\texttt{MOD\_MONTGOMERY\_EXP}.

\subsection{Interface}
\begin{table}[H]
\centering
\caption{Ports de \texttt{RSA}}
\begin{tabular}{llcp{5cm}}
\toprule
\textbf{Port} & \textbf{Dir.} & \textbf{Largeur} & \textbf{Description} \\
\midrule
\texttt{clk}       & in  & 1 & Horloge \\
\texttt{rst}       & in  & 1 & Reset \\
\texttt{start}     & in  & 1 & Lancement \\
\texttt{i\_message} & in  & KEY\_WIDTH & Message $m$ \\
\texttt{i\_exp}     & in  & KEY\_WIDTH & Exposant ($e$ ou $d$) \\
\texttt{i\_N}       & in  & KEY\_WIDTH & Module RSA \\
\texttt{o\_done}    & out & 1 & Résultat prêt \\
\texttt{o\_result}  & out & KEY\_WIDTH & $m^{exp} \bmod N$ \\
\bottomrule
\end{tabular}
\end{table}

\subsection{Machine d'états}
\begin{enumerate}
  \item \textbf{IDLE} : capture des entrées, calcul de $m \bmod N$.
  \item \textbf{REDUCE} : stabilisation.
  \item \textbf{START\_EXP} : impulsion \texttt{start} vers
        \texttt{MOD\_MONTGOMERY\_EXP}.
  \item \textbf{WAIT\_EXP} : attente de \texttt{exp\_done}.
  \item \textbf{LATCH} : capture du résultat.
  \item \textbf{DONE\_STATE} : \texttt{o\_done} = 1.
\end{enumerate}

\subsection{Simulation}
\begin{figure}[H]
\centering
\fbox{\parbox{0.9\textwidth}{\centering\vspace{3cm}
  \textit{[Capture chronogramme RSA : signaux clk, start, i\_message,
  i\_exp, i\_N, o\_done, o\_result]}
\vspace{3cm}}}
\caption{Chronogramme de simulation de \texttt{RSA}}
\label{fig:sim_rsa}
\end{figure}


% ----------------------------------------------------------------------
\section{RSA\_KEYGEN --- Génération de clés}
\label{sec:keygen}

\subsection{Fonction}
Orchestre la génération complète d'une paire de clés RSA :
\begin{enumerate}
  \item Génère un candidat $p$ via \texttt{RANDOM\_GEN}.
  \item Teste $p$ avec \texttt{MILLER\_RABIN}; si composé, retour à 1.
  \item Génère un candidat $q$; teste sa primalité.
  \item Calcule $N = p \times q$ et $\varphi = (p-1)(q-1)$.
  \item Calcule $d = e^{-1} \bmod \varphi$ via \texttt{EXT\_GCD}.
  \item Sort $(N, e, d)$.
\end{enumerate}

\subsection{Interface}
\begin{table}[H]
\centering
\caption{Ports de \texttt{RSA\_KEYGEN}}
\begin{tabular}{llcp{5cm}}
\toprule
\textbf{Port} & \textbf{Dir.} & \textbf{Largeur} & \textbf{Description} \\
\midrule
\texttt{clk}     & in  & 1   & Horloge \\
\texttt{rst}     & in  & 1   & Reset \\
\texttt{start}   & in  & 1   & Lancement génération \\
\texttt{seed}    & in  & 128 & Graine PRNG \\
\texttt{load}    & in  & 1   & Chargement graine \\
\texttt{o\_done}  & out & 1   & Clés prêtes \\
\texttt{o\_N}     & out & KEY\_WIDTH & Module public \\
\texttt{o\_e}     & out & KEY\_WIDTH & Exposant public \\
\texttt{o\_d}     & out & KEY\_WIDTH & Exposant privé \\
\texttt{o\_valid} & out & 1   & \texttt{'1'} si génération réussie \\
\bottomrule
\end{tabular}
\end{table}

\subsection{Machine d'états}
\begin{enumerate}
  \item \textbf{IDLE} : attente.
  \item \textbf{ST\_GEN\_P / ST\_GEN\_P\_WAIT} : génération et attente de $p$.
  \item \textbf{ST\_TEST\_P / ST\_TEST\_P\_WAIT} : test de primalité de $p$.
  \item \textbf{ST\_GEN\_Q / ST\_GEN\_Q\_WAIT} : génération de $q$.
  \item \textbf{ST\_TEST\_Q / ST\_TEST\_Q\_WAIT} : test de primalité de $q$
        (vérifie aussi $p \neq q$).
  \item \textbf{ST\_COMPUTE\_N} : calcul de $N$ et $\varphi$.
  \item \textbf{ST\_GCD\_LAUNCH / ST\_GCD\_WAIT} : calcul de $d$.
  \item \textbf{ST\_OUTPUT\_KEYS} : présentation des résultats.
  \item \textbf{DONE\_ST} : fin.
\end{enumerate}

\subsection{Simulation}
\begin{figure}[H]
\centering
\fbox{\parbox{0.9\textwidth}{\centering\vspace{3cm}
  \textit{[Capture chronogramme RSA\_KEYGEN : signaux clk, start, state,
  o\_done, o\_N, o\_e, o\_d, o\_valid]}
\vspace{3cm}}}
\caption{Chronogramme de simulation de \texttt{RSA\_KEYGEN}}
\label{fig:sim_keygen}
\end{figure}


% ======================================================================
%                  CHAPTER 5 : RSA_TOP (ENTITÉ PRINCIPALE)
% ======================================================================
\chapter{RSA\_TOP --- Entité Principale}
\label{ch:top}

\section{Fonction}

\texttt{RSA\_TOP} est le point d'entrée unique du co-processeur. Il intègre
\texttt{RSA\_KEYGEN} et \texttt{RSA} dans un seul bloc. L'utilisateur
ne fournit que :
\begin{itemize}
  \item \texttt{i\_message} : le message à chiffrer ou déchiffrer.
  \item \texttt{mode} : \texttt{'0'} = chiffrement, \texttt{'1'} = déchiffrement.
  \item \texttt{start} : impulsion de démarrage.
\end{itemize}

Tout le reste (génération de clés, sélection d'exposant, réduction modulaire)
est géré en interne.

\section{Interface externe}
\begin{table}[H]
\centering
\caption{Ports de \texttt{RSA\_TOP}}
\begin{tabular}{llcp{5cm}}
\toprule
\textbf{Port} & \textbf{Dir.} & \textbf{Largeur} & \textbf{Description} \\
\midrule
\texttt{clk}       & in  & 1 & Horloge \\
\texttt{rst}       & in  & 1 & Reset asynchrone \\
\texttt{i\_message} & in  & KEY\_WIDTH & Message (clair ou chiffré) \\
\texttt{mode}      & in  & 1 & '0'=encrypt, '1'=decrypt \\
\texttt{start}     & in  & 1 & Impulsion de démarrage \\
\texttt{o\_done}    & out & 1 & Résultat disponible \\
\texttt{o\_result}  & out & KEY\_WIDTH & Résultat (chiffré ou déchiffré) \\
\texttt{o\_N}       & out & KEY\_WIDTH & Module généré (debug) \\
\texttt{o\_e}       & out & KEY\_WIDTH & Exposant public (debug) \\
\texttt{o\_d}       & out & KEY\_WIDTH & Exposant privé (debug) \\
\bottomrule
\end{tabular}
\end{table}

\section{Séquence d'opérations}

\begin{enumerate}
  \item \textbf{Après reset :} le FSM charge automatiquement la graine
        PRNG et lance \texttt{RSA\_KEYGEN}. Aucune intervention de
        l'utilisateur n'est nécessaire.
  \item \textbf{Keygen terminé :} $N$, $e$, $d$ sont stockés dans des
        registres internes. Le système passe en état \texttt{S\_READY}.
  \item \textbf{Opération RSA :} sur \texttt{start}, le message et le mode
        sont capturés. L'exposant approprié ($e$ ou $d$) est sélectionné.
        Le cœur RSA est lancé.
  \item \textbf{Résultat :} \texttt{o\_done} pulse pendant un cycle
        lorsque \texttt{o\_result} est valide.
\end{enumerate}


\section{Machine d'états principale}

\begin{figure}[H]
\centering
\begin{tikzpicture}[
  state/.style={draw, rounded corners, minimum width=2.5cm,
                minimum height=0.8cm, align=center, font=\small},
  arr/.style={->, thick, >=stealth}
]
\node[state, fill=gray!20]  (seed)   at (0,0)    {S\_SEED\_LOAD};
\node[state, fill=orange!20](kgl)    at (0,-1.5) {S\_KG\_LAUNCH};
\node[state, fill=orange!20](kgr)    at (0,-3)   {S\_KG\_RUN};
\node[state, fill=green!20] (ready)  at (0,-4.5) {S\_READY};
\node[state, fill=blue!15]  (prep)   at (0,-6)   {S\_PREP};
\node[state, fill=blue!15]  (launch) at (0,-7.5) {S\_LAUNCH};
\node[state, fill=blue!15]  (wait)   at (0,-9)   {S\_WAIT\_RSA};
\node[state, fill=blue!15]  (dout)   at (0,-10.5){S\_DONE\_OUT};

\draw[arr] (seed)   -- (kgl);
\draw[arr] (kgl)    -- (kgr);
\draw[arr] (kgr)    -- node[right]{\scriptsize kg\_done \& valid} (ready);
\draw[arr] (kgr.west) -- ++(-1.5,0) |- node[left,pos=0.25]
           {\scriptsize !valid} (kgl.west);
\draw[arr] (ready)  -- node[right]{\scriptsize start} (prep);
\draw[arr] (prep)   -- (launch);
\draw[arr] (launch) -- (wait);
\draw[arr] (wait)   -- node[right]{\scriptsize rsa\_done} (dout);
\draw[arr] (dout.east) -- ++(2,0) |- (ready.east);
\end{tikzpicture}
\caption{Diagramme d'états de \texttt{RSA\_TOP}}
\label{fig:fsm_top}
\end{figure}

\section{Simulation complète (round-trip)}
\begin{figure}[H]
\centering
\fbox{\parbox{0.9\textwidth}{\centering\vspace{4cm}
  \textit{[Capture chronogramme RSA\_TOP : round-trip encrypt+decrypt.
  Signaux : clk, rst, start, i\_message, mode, o\_done, o\_result, o\_N]}
\vspace{4cm}}}
\caption{Simulation round-trip de \texttt{RSA\_TOP} : chiffrement puis déchiffrement}
\label{fig:sim_top}
\end{figure}


% ======================================================================
%                   CHAPTER 6 : RÉSULTATS DE SIMULATION
% ======================================================================
\chapter{Résultats de Simulation}
\label{ch:sim}

\section{Environnement de simulation}
\begin{itemize}
  \item \textbf{Outil :} ModelSim / Vivado Simulator / GHDL
        \textit{[à compléter]}
  \item \textbf{Période d'horloge :} 10~ns (100~MHz)
  \item \textbf{Testbench :} \texttt{TB\_RSA\_TOP.vhd}
\end{itemize}

\section{Scénarios de test}

Le testbench \texttt{TB\_RSA\_TOP} vérifie automatiquement le
\textit{round-trip} : $\text{decrypt}(\text{encrypt}(m)) = m$.

\begin{table}[H]
\centering
\caption{Vecteurs de test du round-trip}
\begin{tabular}{cccl}
\toprule
\textbf{Message $m$} & \textbf{Encrypt} & \textbf{Decrypt} & \textbf{Statut} \\
\midrule
42    & $c = m^e \bmod N$ & $m' = c^d \bmod N$ & \textit{[PASS/FAIL]} \\
99    & $c = m^e \bmod N$ & $m' = c^d \bmod N$ & \textit{[PASS/FAIL]} \\
12345 & $c = m^e \bmod N$ & $m' = c^d \bmod N$ & \textit{[PASS/FAIL]} \\
\bottomrule
\end{tabular}
\end{table}

\section{Performance mesurée}

\begin{table}[H]
\centering
\caption{Nombre de cycles par opération (K = 1024 bits, simulation)}
\begin{tabular}{lr}
\toprule
\textbf{Opération} & \textbf{Cycles approximatifs} \\
\midrule
Génération de clés (keygen)                     & Variable (dépend des candidats) \\
Chiffrement ($e = 65537$, popcount = 2)         & $\sim 40\,000$ \\
Déchiffrement ($d \approx K$ bits, dense)       & $\sim 2\,100\,000$ \\
\textbf{Pire cas} (exposant tous bits à 1)      & $\sim 2 \times K^2 \approx 2\,097\,152$ \\
\bottomrule
\end{tabular}
\end{table}

\textbf{Note :} Avec $e = 65537$ (seulement 2 bits à 1), le chiffrement
est beaucoup plus rapide que le déchiffrement (exposant $d$ de 1024 bits
avec $\sim 50\%$ de bits à 1).

\section{Résultat global de simulation}
\begin{figure}[H]
\centering
\fbox{\parbox{0.9\textwidth}{\centering\vspace{4cm}
  \textit{[Capture d'écran complète de la simulation : fenêtre waveform
  avec tous les signaux principaux, montrant le keygen suivi de 3
  round-trips réussis]}
\vspace{4cm}}}
\caption{Vue d'ensemble de la simulation \texttt{TB\_RSA\_TOP}}
\label{fig:sim_full}
\end{figure}


% ======================================================================
%                 CHAPTER 7 : INTÉGRATION RISC-V
% ======================================================================
\chapter{Intégration avec le Processeur RISC-V}
\label{ch:riscv}

\section{Principe d'intégration}

Le co-processeur RSA est connecté au processeur RISC-V via un bus
mémoire-mappé (AXI-Lite ou Wishbone). Le processeur écrit le message
dans un registre, lance l'opération, et lit le résultat une fois
\texttt{o\_done} actif.

\begin{figure}[H]
\centering
\begin{tikzpicture}[
  block/.style={draw, thick, minimum width=3cm, minimum height=1.5cm,
                align=center, rounded corners=3pt},
  arr/.style={<->, thick, >=stealth}
]
\node[block, fill=blue!15] (rv) at (0,0) {RISC-V\\Core};
\node[block, fill=gray!15] (bus) at (4,0) {Bus\\AXI-Lite};
\node[block, fill=green!15] (rsa) at (8,0) {RSA\_TOP\\Co-proc};

\draw[arr] (rv) -- (bus);
\draw[arr] (bus) -- (rsa);
\end{tikzpicture}
\caption{Intégration du co-processeur RSA avec le RISC-V}
\label{fig:integration}
\end{figure}

\section{Registres mémoire-mappés (proposition)}

\begin{table}[H]
\centering
\caption{Carte mémoire du co-processeur (exemple)}
\begin{tabular}{lcp{6cm}}
\toprule
\textbf{Offset} & \textbf{R/W} & \textbf{Description} \\
\midrule
\texttt{0x00} & R/W & Registre de contrôle (start, mode, reset) \\
\texttt{0x04} & R   & Registre de statut (done, key\_valid) \\
\texttt{0x08--0x87} & W & Message d'entrée (128 octets = 1024 bits) \\
\texttt{0x88--0x107} & R & Résultat (128 octets) \\
\texttt{0x108--0x187} & R & Module $N$ (debug, optionnel) \\
\bottomrule
\end{tabular}
\end{table}

% ======================================================================
%                   CHAPTER 8 : CONCLUSION
% ======================================================================
\chapter{Conclusion et Perspectives}
\label{ch:conclusion}

\section{Bilan}
Ce projet a permis de concevoir et valider un co-processeur cryptographique
RSA complet en VHDL, intégrant :
\begin{itemize}
  \item Un générateur de nombres pseudo-aléatoires (LFSR 128 bits).
  \item Un test de primalité matériel (Miller-Rabin, 4 témoins).
  \item Un calculateur d'inverse modulaire (Euclide étendu séquentiel).
  \item Un cœur d'exponentiation Montgomery bit-sériel.
  \item Une entité top-level autonome (keygen + encrypt/decrypt).
\end{itemize}

La simulation confirme le fonctionnement correct du round-trip
$\text{decrypt}(\text{encrypt}(m)) = m$ pour plusieurs vecteurs de test.


\section{Limites actuelles}
\begin{itemize}
  \item \textbf{Performance :} Le multiplieur bit-sériel traite un seul
        bit par cycle. Un déchiffrement RSA-1024 nécessite environ
        2 millions de cycles ($\sim 20$~ms à 100~MHz).
  \item \textbf{PRNG :} Le LFSR n'est pas cryptographiquement sûr;
        acceptable pour la simulation mais insuffisant pour la production.
  \item \textbf{Multiplication combinatoire :} Les produits $p \times q$
        et $(p-1)(q-1)$ dans \texttt{RSA\_KEYGEN} sont combinatoires
        (non synthétisables pour $K > 64$).
  \item \textbf{Réduction combinatoire :} Le \texttt{mod} dans
        \texttt{RSA.vhd} est combinatoire; nécessite un réducteur
        séquentiel pour la synthèse.
\end{itemize}

\section{Perspectives d'amélioration}
\begin{itemize}
  \item \textbf{Multiplieur radix-4/8 :} traiter 2--3 bits par cycle
        pour une accélération de $2\times$--$3\times$.
  \item \textbf{Double multiplieur :} exécuter carré et multiplication
        en parallèle (gain de $\sim 30\%$).
  \item \textbf{CRT (Chinese Remainder Theorem) :} effectuer deux
        exponentiations de 512 bits au lieu d'une de 1024 bits
        (accélération $\times 4$).
  \item \textbf{Sécurisation :} protection contre les attaques par
        canaux auxiliaires (masquage, randomisation de l'exposant).
  \item \textbf{Intégration bus :} développer le wrapper AXI-Lite
        pour une connexion directe au RISC-V.
  \item \textbf{PRNG robuste :} remplacer le LFSR par un AES-CTR-DRBG.
\end{itemize}

% ======================================================================
%                         ANNEXES
% ======================================================================
\appendix
\chapter{Code VHDL --- Extraits}

\section{MONTGOMERY\_MULT (complet)}
\begin{lstlisting}[style=vhdl, caption={MONTGOMERY\_MULT.vhd}]
architecture rtl of MONTGOMERY_MULT is
  type state_t is (IDLE, RUN, DONE);
  signal state : state_t := IDLE;
  signal X_reg, Y_reg, M_reg : unsigned(K-1 downto 0);
  signal S : unsigned(K+2 downto 0);
  signal cnt : integer range 0 to K-1;
begin
  process(clk, rst)
    variable temp : unsigned(K+2 downto 0);
  begin
    if rst = '1' then
      state <= IDLE; o_done <= '0'; o_Z <= (others => '0');
    elsif rising_edge(clk) then
      o_done <= '0';
      case state is
        when IDLE =>
          if start = '1' then
            X_reg <= i_X; Y_reg <= i_Y; M_reg <= i_M;
            S <= (others => '0'); cnt <= 0; state <= RUN;
          end if;
        when RUN =>
          temp := S;
          if X_reg(cnt) = '1' then temp := temp + Y_reg; end if;
          if temp(0) = '1' then temp := temp + M_reg; end if;
          S <= shift_right(temp, 1);
          if cnt = K-1 then state <= DONE;
          else cnt <= cnt + 1; end if;
        when DONE =>
          if S >= M_reg then o_Z <= resize(S - M_reg, K);
          else o_Z <= resize(S, K); end if;
          o_done <= '1'; state <= IDLE;
      end case;
    end if;
  end process;
end rtl;
\end{lstlisting}


\section{RSA\_TOP --- Machine d'états (extrait)}
\begin{lstlisting}[style=vhdl, caption={FSM principale de RSA\_TOP}]
case state is
  when S_SEED_LOAD =>
    kg_load <= '1';
    state   <= S_KG_LAUNCH;
  when S_KG_LAUNCH =>
    kg_start <= '1';
    state    <= S_KG_RUN;
  when S_KG_RUN =>
    if kg_done = '1' then
      if kg_valid = '1' then
        N_reg <= kg_N; e_reg <= kg_e; d_reg <= kg_d;
        state <= S_READY;
      else
        kg_start <= '1'; -- retry
      end if;
    end if;
  when S_READY =>
    if start = '1' then
      msg_reg  <= i_message;
      mode_reg <= mode;
      state    <= S_PREP;
    end if;
  when S_PREP =>
    rsa_msg <= msg_reg; rsa_N <= N_reg;
    if mode_reg = '1' then rsa_exp <= d_reg;
    else rsa_exp <= e_reg; end if;
    state <= S_LAUNCH;
  when S_LAUNCH =>
    rsa_start <= '1';
    state     <= S_WAIT_RSA;
  when S_WAIT_RSA =>
    if rsa_done = '1' then
      o_result <= rsa_result;
      state    <= S_DONE_OUT;
    end if;
  when S_DONE_OUT =>
    o_done <= '1';
    state  <= S_READY;
end case;
\end{lstlisting}

% ======================================================================
%                       BIBLIOGRAPHIE
% ======================================================================
\begin{thebibliography}{9}

\bibitem{rsa}
R.~Rivest, A.~Shamir, L.~Adleman,
\textit{A Method for Obtaining Digital Signatures and Public-Key Cryptosystems},
Communications of the ACM, vol.~21, no.~2, pp.~120--126, 1978.

\bibitem{montgomery}
P.~L.~Montgomery,
\textit{Modular Multiplication Without Trial Division},
Mathematics of Computation, vol.~44, no.~170, pp.~519--521, 1985.

\bibitem{millerrabin}
M.~O.~Rabin,
\textit{Probabilistic Algorithm for Testing Primality},
Journal of Number Theory, vol.~12, no.~1, pp.~128--138, 1980.

\bibitem{menezes}
A.~J.~Menezes, P.~C.~van Oorschot, S.~A.~Vanstone,
\textit{Handbook of Applied Cryptography},
CRC Press, 1996.

\bibitem{riscv}
A.~Waterman, K.~Asanovic (eds.),
\textit{The RISC-V Instruction Set Manual},
Volume I: Unprivileged ISA, v.~20191213, 2019.

\end{thebibliography}

\end{document}
